\documentclass{article}
\usepackage[UTF8]{ctex}
\usepackage{amsmath,amssymb,amsfonts}
\usepackage{enumitem}
\usepackage{geometry}
\geometry{a4paper, margin=1in}
% === 错题本排版优化（不参与题目标识） ===
\setlength{\parindent}{0pt}
\setlength{\parskip}{0.35em plus 0.10em minus 0.08em}
\setlength{\emergencystretch}{3em}
\setlist{itemsep=0.25em, topsep=0.40em, parsep=0pt}
\allowdisplaybreaks[2]
% === 错题本排版优化结束 ===
\begin{document}
\section*{华东理工大学2024年高等代数真题}
\section*{46．华东理工大学2024年高等代数试卷（回忆版）}

\subsection{题目1}
一．求一个三次多项式 \(f(x)\) 使得 \((x-1)^{2}\left|f(x)+1,(x+1)^{2}\right| f(x)-1\) ．

\subsection{题目2}
二．设矩阵

\[
A=\left(\begin{array}{lll}
1 & 0 & 0 \\
1 & 1 & 0 \\
1 & 1 & 1
\end{array}\right), B=\left(\begin{array}{lll}
0 & 1 & 1 \\
1 & 0 & 1 \\
1 & 1 & 0
\end{array}\right) .
\]

若矩阵 \(X\) 满足 \(A X A+B X B=A X B+B X A+E\) ，求 \(X\) 。

\subsection{题目3}
三．设行列式

\[
D=\left|\begin{array}{cccc}
a_{11} & a_{12} & \cdots & a_{1 n} \\
\vdots & \vdots & & \vdots \\
a_{n-1,1} & a_{n-1,2} & \cdots & a_{n-1, n} \\
1 & 1 & \cdots & 1
\end{array}\right|
\]

\(D_{i}\) 是用 \(\left(x_{1}, x_{2}, \cdots, x_{n-1}, 1\right)^{\mathrm{T}}\) 替换掉 \(D\) 的第 \(i\) 列后得到的行列式，证明：

\[
D_{1}+D_{2}+\cdots+D_{n}=D .
\]

\subsection{题目4}
四．设 \(A\) 是秩为 \(r\) 的 \(m \times n\) 矩阵，\(B\) 是秩为 \(n-r\) 的 \(n \times p\) 矩阵，且 \(A B=O\) ．若向量 \(b\) 满足 \(A b=0\)－证明：

\begin{itemize}
  \item[（1）] 方程组 \(B X=b\) 有解。
  \item[（2）] 当 \(p=n-r\) 时，方程组 \(B X=b\) 有唯一解。
\end{itemize}

\subsection{题目5}
五。设 \(\alpha_{1}, \alpha_{2}, \cdots, \alpha_{n}\) 为 \(n\) 个列向量，满足

\[
\alpha_{i}^{\mathrm{T}} \alpha_{j}= \pm \cos \theta, i \neq j, \cos \theta \neq 1 .
\]

记 \(X_{i}=\alpha_{i} \alpha_{i}^{\mathrm{T}}, i=1,2, \cdots, n\) 。

\begin{itemize}
  \item[（1）] 求 \(\operatorname{tr}\left(X_{i}\right)\) 和 \(\operatorname{tr}\left(X_{i} X_{j}\right)\) 。
  \item[（2）] 证明：\(X_{1}, X_{2}, \cdots, X_{n}\) 线性无关。
\end{itemize}

\subsection{题目6}
六。设实矩阵 \(A\) 满足 \(A+A^{\mathrm{T}}=J-E\) ，其中 \(J\) 是元素全为1的 \(n\) 阶方阵。证明：

\begin{itemize}
  \item[（1）] 对任意复向量 \(\alpha\) ，有 \(0 \leq \alpha^{H} J \alpha \leq n \alpha^{H} \alpha\) 。
  \item[（2）] \(A\) 的特征值在 \(\left[-\frac{1}{2}, \frac{n-1}{2}\right]\) 之间．
\end{itemize}

\subsection{题目7}
七．在所有 \(n\) 阶实对称矩阵构成的线性空间 \(V\) 上定义内积 \((A, B)=\operatorname{tr}(A B)\) ．

\begin{itemize}
  \item[（1）] 证明：\(V\) 构成欧几里得空间。
  \item[（2）] \(S\) 是 \(V\) 中满足 \(\operatorname{tr}(A)=0\) 的矩阵组成的子空间，求 \(\operatorname{dim} S\) ．
  \item[（3）] 求 \(\operatorname{dim}\left(S^{\perp}\right)\) ．
\end{itemize}

\subsection{题目8}
八．设 \(V\) 是西空间，\(W\) 是 \(V\) 的子空间，\(\beta \in V\) 是一个向量，如果 \(W\) 中的向量 \(\alpha\) 满足对任意的 \(\gamma \in W\)都有 \(\|\beta-\alpha\| \leq\|\beta-\gamma\|\) ，则称 \(\alpha\) 是 \(\beta\) 到 \(W\) 的投影。

\begin{itemize}
  \item[（1）] \(\alpha\) 属于 \(W\) 的正交补。
  \item[（2）] \(\alpha\) 是唯一的．
\end{itemize}

\subsection{题目9}
九．设 \(n\) 维欧氏空间 \(V\) 上的两个线性变换 \(\sigma, \tau\) 满足

\[
(\sigma(\alpha), \sigma(\beta))=(\tau(\alpha), \tau(\beta)), \forall \alpha, \beta \in V .
\]

证明： \(\operatorname{Im} \sigma\) 和 \(\operatorname{Im} \tau\) 同构．

\subsection{题目10}
十．线性空间 \(V\) 和 \(V\) 的子空间 \(V_{1}, V_{2}\) 满足 \(V=V_{1} \oplus V_{2}, \sigma, \tau\) 是两个线性变换，\(\sigma(\alpha)=\alpha_{1}+\alpha_{2}\left(\alpha_{i} \in V_{i}\right)\)证明：\(V_{1}, V_{2}\) 都是 \(\tau\) 的不变子空间的充要条件是 \(\sigma \tau=\tau \sigma\) ．
\end{document}
